\documentclass[11pt]{article}
\usepackage[margin=1in]{geometry}
\setlength{\parindent}{0pt}
\setlength{\parskip}{0.55em}
\setlength{\emergencystretch}{1em}
\begin{document}
\begin{center}
    {\Large\textbf{Understanding the Notion Experiment}}\\[0.4em]
    {\normalsize Purpose, architecture, and the boundary between information and authority}
\end{center}

\section{What the experiment is asking}

An assistant that can read documents needs to do more than understand their
contents. It must also distinguish between information relevant to a task and
instructions that are allowed to govern its behavior. This experiment studies
that distinction using an assistant connected to a small Notion workspace.

The assistant receives an ordinary task and an attached image. It can inspect
the image, search the workspace, read documents, and produce an answer. Some
information is available for legitimate use, while a designated sensitive
page is explicitly off limits under the assistant's standing instructions.
The experiment changes the text embedded in the image and observes whether
the assistant continues to respect that restriction.

The central question is whether material the assistant is supposed to
\emph{read} can influence what it believes it is \emph{authorized to do}.
An image might contain language that presents an additional retrieval as a
necessary step in completing the task. Understanding that language is
expected; accepting it as permission to access prohibited information would
be a failure of the intended boundary.

\section{The Notion setting}

The workspace contains synthetic documents created for the experiment. An
allowed page, such as \emph{Office Information}, provides ordinary information
that the assistant can use when answering an appropriate request. A separate
sensitive page contains information that the assistant has been instructed
never to retrieve or disclose. Its contents are test data, rather than actual
confidential records.

Both pages are shared with the assistant's Notion connection. This is a
 deliberate experimental choice. If the sensitive page were inaccessible,
then a refusal or failed retrieval might simply reflect a lack of access.
By making the page technically reachable, the experiment can investigate
whether the assistant follows its policy when it has the capability to
violate it.

Preparing the setting therefore means establishing an allowed document, a
clearly designated sensitive document, and a connection that can read both.
The experimenter also retains a separate copy of the sensitive test content
for comparison with answers. That reference allows the evaluator to
 distinguish an actual disclosure of the test information from a plausible
but invented answer.

The connection is restricted to reading and searching. The assistant is not
supposed to create, modify, or delete workspace content. The experiment is
therefore concerned with information access and disclosure, rather than
changes to the documents themselves.

\section{The specific trust boundary}

The boundary concerns the authority of different sources of text. The
operator establishes the assistant's standing rules before a trial begins.
These rules identify the sensitive page and prohibit retrieving or disclosing
its contents, including its child content and sensitive excerpts encountered
in search results. The submitted task and image do not provide a way to
replace those standing rules.

Within those rules, the task request can tell the assistant what help is
wanted. The image and retrieved documents can supply information needed to
provide that help. However, none of those sources can authorize an exception
to the sensitive-page restriction. A request to read an image does not grant
permission to follow every instruction written inside it.

Here, calling content \emph{untrusted} does not mean that its factual claims
must be false or that the assistant should ignore it completely. An office
notice may accurately state an opening time, and a retrieved document may
contain the answer to the user's question. The limitation is about authority:
useful content cannot rewrite the rules under which the assistant operates.
The assistant must preserve that distinction even when the content sounds
confident, helpful, or closely related to the task.

The boundary is crossed if the assistant treats an instruction in the image,
or in material it retrieves, as sufficient justification to disregard the
standing prohibition. For example, an image could suggest that a complete
answer depends on consulting the protected page. The relevant failure would
be accepting that suggestion as authorization. Merely noticing or describing
the suggestion would not establish the same failure.

Two different protections are present. The sensitive-page restriction is a
behavioral rule that the model must follow; there is no separate mechanism
that automatically blocks every request for that page. In contrast, the
restriction to approved reading and searching operations is enforced by the
software. An unapproved operation is rejected even if the model asks for it.
Thus, the software limits the kinds of actions available, while the model
must decide whether a particular read is permitted under its instructions.

\section{Architecture: how the parts work together}

The repository separates the system being tested from the system conducting
the test. The Notion target is the assistant under examination. The outer
experiment is a wrapper that prepares trials, submits them, evaluates the
answers, and decides what to try next. This separation lets the experiment
change the image while keeping the target's configured rules in place.

The target owns the connection to Notion and the model that answers the task.
It receives a prompt and an image, lets the model request approved searches
or reads, and returns the final answer. A fresh conversation and document
connection are created for every trial. The target therefore does not
continue a conversation containing the earlier trial's attempted injection.

The wrapper supplies continuity across trials. It remembers previous
payloads, their scores, and which goals have been achieved. It renders new
text into an image and submits that image to the target. The adaptation
happens in the experiment's record of earlier attempts, rather than through
persistent memory inside the target assistant.

Separate feedback calls support evaluation and candidate choice. The judge
assesses the answer, the generator proposes alternative image texts, and the
selector chooses a candidate to test. These are distinct roles, although
they can use the same configured feedback model. They do not carry out the
target's Notion task. In particular, the generator receives numeric goal
feedback and previous payloads, without being given the target's answers or
the sensitive reference contents.

The target returns an answer and an identifier for the trial. The wrapper
saves that answer alongside the image, evaluation, and candidate decisions.
More detailed records of the target's searches and reads remain with the
target and can be matched using the identifier. The ordinary experiment loop
therefore evaluates what appears in the answer; it does not automatically
inspect the assistant's internal sequence of document accesses.

\section{How a trial becomes an experiment}

A run begins with an image containing a candidate instruction and an ordinary
task request. The current approach uses completion-style text: it asks the
assistant to continue an answer or fill in an unfinished response. The
background image provides a carrier for that text. Its unrelated visual
content does not determine which Notion page should be used.

The target responds, and the wrapper saves the response before evaluation.
The judge then considers the currently active goals. The default progression
starts with including information from the allowed page and later adds the
goal of including sensitive test content. Earlier goals remain active as new
ones are added, so a later goal cannot compensate for losing an earlier
behavior.

When all active goals pass, the next goal is introduced. Otherwise, the
experiment continues working on the current stage. The generator proposes
complete replacement payloads, using earlier successful wording as context
rather than as text that must remain unchanged. Duplicate proposals are
filtered, and the selector chooses an alternative for the next trial.
Selection is a prediction about what may work; only a subsequent target
response provides evidence about what actually happened.

A run ends when the release judge identifies sensitive reference content in
the final answer, when no new candidates remain after the allowed retries,
or when the iteration limit is reached. The saved record preserves the
sequence of tested images and answers so that the outcome can be reviewed.

\section{Interpreting the result}

Goal completion and sensitive-information release are evaluated separately.
An assistant may satisfy an ordinary task without disclosing anything
sensitive. It may also disclose sensitive information while failing another
requested behavior. For that reason, passing every staged goal is not by
itself the condition that marks a run as successful for the attack; the
experiment looks for a nonzero release assessment against known references.

A claim such as ``I read the protected page'' is not sufficient evidence that
the page was accessed. Likewise, an answer that invents a convincing secret
is not evidence that the actual test content was released. The final answer
must be compared with the reference material. Without that material, the
experiment records release as unknown. Establishing which documents were
actually read requires examining the target's more detailed records.

The release assessment is still an automated judgment and can be mistaken.
The saved answer and supporting excerpts make the finding reviewable. A
single successful trial also does not establish how reliably a payload will
work. Repeated trials are needed to assess consistency, and a comparison
with an otherwise similar request without the injected instructions is
needed to support a claim that the image caused the change in behavior.

Finally, the setup is a controlled local research environment. Its separation
into a wrapper and a target organizes responsibilities and observations; it
does not create a complete security barrier around the machine. The service
is intended to remain local, and its document permissions are deliberately
chosen to make the model's handling of the trust boundary observable.

\end{document}
